% GENERATED FILE. Edit canonical lessons, then run npm run generate:books.
% canonical-source-hash: fnv1a64:a16297ce4b317ed8
% canonical-lessons: HI-C477-hatheli, HI-C477-mutthi, HI-C477-bah, HI-C477-tag, HI-C477-eri

\chapter{Arms and Legs}
\label{ch:hi-arms-and-legs}
\hlchaptermodality{477}

\begin{chapteropening}
\textbf{By the end of this chapter:} \emph{I can name parts of the arms and legs.}

Complete the last lesson of chapter 477: I can name parts of the arms and legs.
\end{chapteropening}

\section[\texorpdfstring{\hi{हथेली} (hathelī)}{हथेली (hathelī)}]{\hi{हथेली} (hathelī) — the palm of the hand}

\label{lesson:HI-C477-hatheli}



\begin{quote}
Before the new one: say the Hindi for a beard, then the Hindi for a jaw.
\end{quote}

\subsection*{You'll want to know: \hi{हथेली}}
\textbf{\hi{हथेली}} — \emph{hathelī} — \textquotedblleft{}the palm of the hand\textquotedblright{}.

\begin{quote}
\hi{बच्चे} \hi{की} \hi{हथेली} \hi{बहुत} \hi{छोटी} \hi{और} \hi{नरम} \hi{है।}
\end{quote}

\begin{quote}
\emph{bacce kī hathelī bahut choṭī aur naram hai.}
\end{quote}

\begin{quote}
The baby's palm is very small and soft.
\end{quote}

Feminine.

\begin{cousinweb}
\textbf{\hi{हथेली}} (\emph{hathelī}) is a tadbhava: CDIAL lists Hindi \emph{hathelī}, \textquotedblleft{}palm\textquotedblright{}, under Sanskrit \emph{hastatala}, \textquotedblleft{}the flat of the hand\textquotedblright{} (CDIAL 14029), built on \emph{hasta}, the word that gave \emph{hāth}.
\end{cousinweb}

\begin{grammarlens}[title={What you've built}]
One more word for describing what you lived and felt.
\end{grammarlens}

\subsection*{Your turn}
\begin{itemize}
\raggedright
  \item \emph{Say it:} \emph{hathelī}
  \item \emph{Say it:} \emph{hathelī}, once more
  \item \emph{Say it:} \emph{bacce kī hathelī bahut choṭī aur naram hai.} — The baby's palm is very small and soft.
  \item \emph{Recall:} say the Hindi for a beard, then the Hindi for a jaw, then say \emph{hathelī} again
\end{itemize}

\begin{tcolorbox}[breakable,colback=teal!4,colframe=teal!35!black,title={Before you move on}]
What does \hi{हथेली} mean? (\textquotedblleft{}The palm of the hand\textquotedblright{}.) Say it once more.
\end{tcolorbox}
\section[\texorpdfstring{\hi{मुट्ठी} (muṭṭhī)}{मुट्ठी (muṭṭhī)}]{\hi{मुट्ठी} (muṭṭhī) — a fist; a handful}

\label{lesson:HI-C477-mutthi}



\begin{quote}
Before the new one: say the Hindi for a jaw, then the Hindi for the palm of the hand.
\end{quote}

\subsection*{You'll want to know: \hi{मुट्ठी}}
\textbf{\hi{मुट्ठी}} — \emph{muṭṭhī} — \textquotedblleft{}a fist; a handful\textquotedblright{}.

\begin{quote}
\hi{बच्चे} \hi{ने} \hi{मुट्ठी} \hi{में} \hi{एक} \hi{सिक्का} \hi{छिपाया।}
\end{quote}

\begin{quote}
\emph{bacce ne muṭṭhī meṁ ek sikkā chipāyā.}
\end{quote}

\begin{quote}
The child hid a coin in his fist.
\end{quote}

Feminine.

\begin{cousinweb}
\textbf{\hi{मुट्ठी}} (\emph{muṭṭhī}) is a tadbhava: CDIAL lists Hindi \emph{muṭṭhī}, \textquotedblleft{}fist, handful\textquotedblright{}, under Sanskrit \emph{muṣṭí}, \textquotedblleft{}fist\textquotedblright{} (CDIAL 10221).
\end{cousinweb}

\begin{grammarlens}[title={What you've built}]
One more word for describing what you lived and felt.
\end{grammarlens}

\subsection*{Your turn}
\begin{itemize}
\raggedright
  \item \emph{Say it:} \emph{muṭṭhī}
  \item \emph{Say it:} \emph{muṭṭhī}, once more
  \item \emph{Say it:} \emph{bacce ne muṭṭhī meṁ ek sikkā chipāyā.} — The child hid a coin in his fist.
  \item \emph{Recall:} say the Hindi for a jaw, then the Hindi for the palm of the hand, then say \emph{muṭṭhī} again
\end{itemize}

\begin{tcolorbox}[breakable,colback=teal!4,colframe=teal!35!black,title={Before you move on}]
What does \hi{मुट्ठी} mean? (\textquotedblleft{}A fist; a handful\textquotedblright{}.) Say it once more.
\end{tcolorbox}
\section[\texorpdfstring{\hi{बाँह} (bā̃h)}{बाँह (bā̃h)}]{\hi{बाँह} (bā̃h) — an arm}

\label{lesson:HI-C477-bah}



\begin{quote}
Before the new one: say the Hindi for the palm of the hand, then the Hindi for a fist.
\end{quote}

\subsection*{You'll want to know: \hi{बाँह}}
\textbf{\hi{बाँह}} — \emph{bā̃h} — \textquotedblleft{}an arm\textquotedblright{}.

\begin{quote}
\hi{माँ} \hi{ने} \hi{बच्चे} \hi{को} \hi{बाँहों} \hi{में} \hi{उठाया।}
\end{quote}

\begin{quote}
\emph{māṁ ne bacce ko bā̃hoṁ meṁ uṭhāyā.}
\end{quote}

\begin{quote}
The mother lifted the child in her arms.
\end{quote}

Feminine.

\begin{cousinweb}
\textbf{\hi{बाँह}} (\emph{bā̃h}) is a tadbhava: CDIAL lists Hindi \emph{bā̃h}, \textquotedblleft{}upper arm\textquotedblright{}, under Sanskrit \emph{bāhú}, \textquotedblleft{}arm\textquotedblright{} (CDIAL 9229).
\end{cousinweb}

\begin{grammarlens}[title={What you've built}]
One more word for describing what you lived and felt.
\end{grammarlens}

\subsection*{Your turn}
\begin{itemize}
\raggedright
  \item \emph{Say it:} \emph{bā̃h}
  \item \emph{Say it:} \emph{bā̃h}, once more
  \item \emph{Say it:} \emph{māṁ ne bacce ko bā̃hoṁ meṁ uṭhāyā.} — The mother lifted the child in her arms.
  \item \emph{Recall:} say the Hindi for the palm of the hand, then the Hindi for a fist, then say \emph{bā̃h} again
\end{itemize}

\begin{tcolorbox}[breakable,colback=teal!4,colframe=teal!35!black,title={Before you move on}]
What does \hi{बाँह} mean? (\textquotedblleft{}An arm\textquotedblright{}.) Say it once more.
\end{tcolorbox}
\section[\texorpdfstring{\hi{टाँग} (ṭā̃g)}{टाँग (ṭā̃g)}]{\hi{टाँग} (ṭā̃g) — a leg}

\label{lesson:HI-C477-tag}



\begin{quote}
Before the new one: say the Hindi for a fist, then the Hindi for an arm.
\end{quote}

\subsection*{You'll want to know: \hi{टाँग}}
\textbf{\hi{टाँग}} — \emph{ṭā̃g} — \textquotedblleft{}a leg\textquotedblright{}.

\begin{quote}
\hi{सीढ़ी} \hi{से} \hi{गिरकर} \hi{अरुण} \hi{की} \hi{टाँग} \hi{टूट} \hi{गई।}
\end{quote}

\begin{quote}
\emph{sīṛhī se girkar Aruṇ kī ṭā̃g ṭūṭ gaī.}
\end{quote}

\begin{quote}
Arun broke his leg falling down the stairs.
\end{quote}

Feminine.

\begin{cousinweb}
\textbf{\hi{टाँग}} (\emph{ṭā̃g}) is inherited: CDIAL lists Hindi \emph{ṭā̃g}, \textquotedblleft{}leg\textquotedblright{}, under Sanskrit \emph{ṭaṅka} (CDIAL 5428). \emph{pair} is the foot; \emph{ṭā̃g} is the whole leg.
\end{cousinweb}

\begin{grammarlens}[title={What you've built}]
One more word for describing what you lived and felt.
\end{grammarlens}

\subsection*{Your turn}
\begin{itemize}
\raggedright
  \item \emph{Say it:} \emph{ṭā̃g}
  \item \emph{Say it:} \emph{ṭā̃g}, once more
  \item \emph{Say it:} \emph{sīṛhī se girkar Aruṇ kī ṭā̃g ṭūṭ gaī.} — Arun broke his leg falling down the stairs.
  \item \emph{Recall:} say the Hindi for a fist, then the Hindi for an arm, then say \emph{ṭā̃g} again
\end{itemize}

\begin{tcolorbox}[breakable,colback=teal!4,colframe=teal!35!black,title={Before you move on}]
What does \hi{टाँग} mean? (\textquotedblleft{}A leg\textquotedblright{}.) Say it once more.
\end{tcolorbox}
\section[\texorpdfstring{\hi{एड़ी} (eṛī)}{एड़ी (eṛī)}]{\hi{एड़ी} (eṛī) — a heel}

\label{lesson:HI-C477-eri}



\begin{quote}
Before the new one: say the Hindi for an arm, then the Hindi for a leg.
\end{quote}

\subsection*{You'll want to know: \hi{एड़ी}}
\textbf{\hi{एड़ी}} — \emph{eṛī} — \textquotedblleft{}a heel\textquotedblright{}.

\begin{quote}
\hi{नए} \hi{जूतों} \hi{से} \hi{मेरी} \hi{एड़ी} \hi{में} \hi{दर्द} \hi{है।}
\end{quote}

\begin{quote}
\emph{nae jūtoṁ se merī eṛī meṁ dard hai.}
\end{quote}

\begin{quote}
My heel hurts because of my new shoes.
\end{quote}

Feminine.

\begin{cousinweb}
\textbf{\hi{एड़ी}} (\emph{eṛī}) is a tadbhava: CDIAL lists Hindi \emph{eṛī}, \textquotedblleft{}heel\textquotedblright{}, under \emph{*aḍḍi}, a form it reconstructs (CDIAL 191).
\end{cousinweb}

\begin{grammarlens}[title={What you've built}]
One more word for describing what you lived and felt.
\end{grammarlens}

\subsection*{Your turn}
\begin{itemize}
\raggedright
  \item \emph{Say it:} \emph{eṛī}
  \item \emph{Say it:} \emph{eṛī}, once more
  \item \emph{Say it:} \emph{nae jūtoṁ se merī eṛī meṁ dard hai.} — My heel hurts because of my new shoes.
  \item \emph{Recall:} say the Hindi for an arm, then the Hindi for a leg, then say \emph{eṛī} again
\end{itemize}

\begin{tcolorbox}[breakable,colback=teal!4,colframe=teal!35!black,title={Before you move on}]
What does \hi{एड़ी} mean? (\textquotedblleft{}A heel\textquotedblright{}.) Say it once more.
\end{tcolorbox}
